\section{Kết quả và insight từ biểu đồ}

\subsection{Đọc insight theo ngôn ngữ trực quan}

Phần này được trình bày theo đúng cách đọc biểu đồ trong Data Visualization:
đầu tiên đọc \textbf{hình dạng tổng thể}, sau đó đọc \textbf{khoảng cách tương đối}, cuối cùng mới gắn với \textbf{kết luận định lượng}.

\subsection{Hình \ref{fig:time-linear}: góc nhìn chênh lệch tuyệt đối}

Trên trục $y$ tuyến tính, điều nhìn thấy rõ nhất là "độ cao tuyệt đối" của thời gian chạy:
\begin{itemize}[leftmargin=1.2cm]
    \item Đường Insertion Sort vọt mạnh ở mốc $n=100000$ trong mọi panel, tạo khoảng cách rất lớn với hai thuật toán còn lại.
    \item Quick Sort và Counting Sort nằm sát đáy đồ thị ở panel có Insertion tăng mạnh, cho thấy hiện tượng "nén thị giác" khi dải giá trị quá rộng.
    \item Vì vậy, linear scale phù hợp để trả lời câu hỏi: \textit{mức tốn thời gian tuyệt đối là bao nhiêu}. Hình này cho thấy rõ chi phí thực thi thực tế nếu ưu tiên thời gian chạy.
\end{itemize}

\subsection{Hình \ref{fig:time-logy}: góc nhìn cấu trúc tăng trưởng}

Khi chuyển sang log-$y$, dữ liệu không còn bị nén mạnh và hình thái đường trở nên dễ so sánh:
\begin{itemize}[leftmargin=1.2cm]
    \item Thứ hạng giữa ba thuật toán ổn định trên tất cả panel: Counting Sort nhanh nhất, Quick Sort đứng thứ hai, Insertion Sort chậm nhất.
    \item Độ dốc tương đối của các đường phản ánh tốc độ tăng khi $n$ tăng: Insertion tăng nhanh nhất, Counting tăng chậm nhất.
    \item Panel Nearly Sorted cho thấy Insertion "dịu" hơn rõ rệt so với Random/Reversed, nhưng vẫn không vượt qua Quick hoặc Counting ở mốc lớn.
\end{itemize}

Nói cách khác, nếu Hình \ref{fig:time-linear} trả lời "\textit{chậm bao nhiêu ms}", thì Hình \ref{fig:time-logy} trả lời "\textit{tăng nhanh đến mức nào theo bậc độ lớn}".

\subsection{Hình \ref{fig:bar-nmax}: snapshot ra quyết định ở tải lớn}

Biểu đồ cột tại $n=100000$ đóng vai trò tổng kết nhanh:
\begin{itemize}[leftmargin=1.2cm]
    \item Trong cả 4 input type, cột Counting Sort luôn thấp nhất, cột Quick Sort luôn ở giữa, cột Insertion Sort luôn cao nhất.
    \item Khoảng cách giữa Insertion và hai thuật toán còn lại rất lớn ở Random, Reversed và Many Duplicates; nhỏ nhất ở Nearly Sorted nhưng vẫn vượt trội theo hướng bất lợi cho Insertion.
    \item Với góc nhìn triển khai thực tế ở dữ liệu lớn, ranking thuật toán là nhất quán và không đảo chiều theo loại dữ liệu.
\end{itemize}

\subsection{Hình \ref{fig:comparisons}: biểu đồ phụ trợ kiểm chứng}

Biểu đồ comparisons trên thang log-log có giá trị như một lớp kiểm chứng:
\begin{itemize}[leftmargin=1.2cm]
    \item Thứ hạng theo số phép so sánh đồng pha với thứ hạng theo thời gian chạy.
    \item Insertion có mức comparisons cao nhất, phù hợp với việc đường thời gian tăng mạnh trong hai hình chính.
    \item Counting có mức comparisons thấp nhất và ổn định hơn, nhất quán với vai trò thuật toán nhanh nhất trên các hình thời gian.
\end{itemize}

Do đó, insight không chỉ dựa vào một chỉ số duy nhất mà được củng cố bởi hai lớp biểu diễn khác nhau: \textit{time} và \textit{comparisons}.

\subsection{Insight theo từng kiểu dữ liệu}

\textbf{Random:}
\begin{itemize}[leftmargin=1.2cm]
    \item Trên cả linear và log-$y$, khoảng cách giữa Insertion và hai thuật toán còn lại mở rộng rất nhanh khi $n$ tăng.
    \item Đây là kịch bản tổng quát cho thấy ưu thế ổn định của Quick và Counting so với Insertion ở dữ liệu lớn.
\end{itemize}

\textbf{Reversed:}
\begin{itemize}[leftmargin=1.2cm]
    \item Panel Reversed thể hiện Insertion tệ nhất trong bốn loại input (đường cao nhất và dốc mạnh nhất ở vùng $n$ lớn).
    \item Quick và Counting duy trì khoảng cách thấp hơn nhiều, cho thấy độ bền tốt hơn trước dữ liệu bất lợi.
\end{itemize}

\textbf{Nearly Sorted:}
\begin{itemize}[leftmargin=1.2cm]
    \item Đây là panel duy nhất mà đường Insertion thấp hơn rõ so với chính nó ở Random/Reversed.
    \item Tuy nhiên về thứ hạng tổng thể, Insertion vẫn đứng sau Quick và Counting ở mốc lớn.
\end{itemize}

\textbf{Many Duplicates:}
\begin{itemize}[leftmargin=1.2cm]
    \item Counting thể hiện ưu thế rõ và ổn định, phù hợp với bản chất thuật toán khi miền giá trị có nhiều phần tử lặp.
    \item Quick vẫn giữ vị trí trung gian; Insertion tăng mạnh ở mốc $n$ lớn.
\end{itemize}

\subsection{Kết luận rút ra trực tiếp từ bộ hình}

\begin{table}[H]
\centering
\renewcommand{\arraystretch}{1.25}
\begin{tabular}{p{3.8cm}p{9.2cm}}
\toprule
\textbf{Khía cạnh} & \textbf{Insight trực quan chính} \\
\midrule
Thứ hạng thuật toán & Trên mọi hình chính, thứ hạng không đảo chiều: Counting Sort $<$ Quick Sort $<$ Insertion Sort. \\
Ảnh hưởng kiểu dữ liệu & Reversed là trường hợp bất lợi nhất cho Insertion; Nearly Sorted là trường hợp Insertion cải thiện rõ nhất. \\
Vai trò của scale & Linear cho đọc chênh lệch tuyệt đối, log-$y$ cho đọc cấu trúc tăng trưởng; cần dùng cả hai để tránh kết luận thiếu. \\
Độ tin cậy insight & Biểu đồ comparisons đồng pha với biểu đồ thời gian, củng cố rằng kết luận không do hiệu ứng hiển thị. \\
\bottomrule
\end{tabular}
\caption{Tổng hợp insight theo góc nhìn trực quan hóa dữ liệu}
\label{tab:top-insights}
\end{table}
